% GENERATED FILE. Edit canonical lessons, then run npm run generate:books.
% canonical-source-hash: fnv1a64:b5816a614e8363fa
% canonical-lessons: UR-C316-qaidi, UR-C316-pahra, UR-C316-surang, UR-C316-khisakna, UR-C316-rengna

\chapter{Escape from the Fort}
\label{ch:ur-escape-from-the-fort}
\hlchaptermodality{316}

\begin{chapteropening}
\textbf{By the end of this chapter:} \emph{I can tell how prisoners got away from a fort.}

Complete the last lesson of chapter 316: I can tell how prisoners got away from a fort.
\end{chapteropening}

\section[\texorpdfstring{\ur{قیدی} (qaidī)}{قیدی (qaidī)}]{\ur{قیدی} (qaidī) — a prisoner}

\label{lesson:UR-C316-qaidi}



\begin{quote}
Before the new one: say the Urdu for a century, then the Urdu for to become deserted.
\end{quote}

\subsection*{You'll want to know: \ur{قیدی}}
\textbf{\ur{قیدی}} — \emph{qaidī} — \textquotedblleft{}a prisoner\textquotedblright{}.

\begin{quote}
\ur{قلعے} \ur{میں} \ur{سو} \ur{قیدی} \ur{تھے۔}
\end{quote}

\begin{quote}
\emph{qile meṅ sau qaidī the.}
\end{quote}

\begin{quote}
There were a hundred prisoners in the fort.
\end{quote}

Masculine or feminine, by the person meant.

\begin{cousinweb}
\textbf{\ur{قیدی}} (\emph{qaidī}) is usually traced, through Persian, to Arabic \emph{qayd}, \textquotedblleft{}a fetter\textquotedblright{}, with the ending \emph{-ī}: one held in chains.
\end{cousinweb}

\begin{grammarlens}[title={What you've built}]
One more word for telling what happened.
\end{grammarlens}

\subsection*{Your turn}
\begin{itemize}
\raggedright
  \item \emph{Say it:} \emph{qaidī}
  \item \emph{Say it:} \emph{qaidī}, once more
  \item \emph{Say it:} \emph{qile meṅ sau qaidī the.} — There were a hundred prisoners in the fort.
  \item \emph{Recall:} say the Urdu for a century, then the Urdu for to become deserted, then say \emph{qaidī} again
\end{itemize}

\begin{tcolorbox}[breakable,colback=teal!4,colframe=teal!35!black,title={Before you move on}]
What does \ur{قیدی} mean? (\textquotedblleft{}A prisoner\textquotedblright{}.) Say it once more.
\end{tcolorbox}
\section[\texorpdfstring{\ur{پہرہ} (pahrā)}{پہرہ (pahrā)}]{\ur{پہرہ} (pahrā) — a guard kept; a watch}

\label{lesson:UR-C316-pahra}



\begin{quote}
Before the new one: say the Urdu for to become deserted, then the Urdu for a prisoner.
\end{quote}

\subsection*{You'll want to know: \ur{پہرہ}}
\textbf{\ur{پہرہ}} — \emph{pahrā} — \textquotedblleft{}a guard kept; a watch\textquotedblright{}.

\begin{quote}
\ur{رات} \ur{کو} \ur{دروازے} \ur{پر} \ur{پہرہ} \ur{تھا۔}
\end{quote}

\begin{quote}
\emph{rāt ko darvāze par pahrā thā.}
\end{quote}

\begin{quote}
At night the gate was guarded.
\end{quote}

Masculine.

\begin{cousinweb}
\textbf{\ur{پہرہ}} (\emph{pahrā}) is usually traced to Persian \emph{pahra}, \textquotedblleft{}a watch, a guard\textquotedblright{}. CDIAL lists the Hindi cousin \emph{paharā}, \textquotedblleft{}a watch, a watchman\textquotedblright{}, under Sanskrit \emph{prahara}, a watch of about three hours, and asks whether that Sanskrit word is itself the Persian one (CDIAL 8900). \textbf{\ur{پہرہ} \ur{دینا}} is to stand guard.
\end{cousinweb}

\begin{grammarlens}[title={What you've built}]
One more word for telling what happened.
\end{grammarlens}

\subsection*{Your turn}
\begin{itemize}
\raggedright
  \item \emph{Say it:} \emph{pahrā}
  \item \emph{Say it:} \emph{pahrā}, once more
  \item \emph{Say it:} \emph{rāt ko darvāze par pahrā thā.} — At night the gate was guarded.
  \item \emph{Recall:} say the Urdu for to become deserted, then the Urdu for a prisoner, then say \emph{pahrā} again
\end{itemize}

\begin{tcolorbox}[breakable,colback=teal!4,colframe=teal!35!black,title={Before you move on}]
What does \ur{پہرہ} mean? (\textquotedblleft{}A guard kept; a watch\textquotedblright{}.) Say it once more.
\end{tcolorbox}
\section[\texorpdfstring{\ur{سرنگ} (surang)}{سرنگ (surang)}]{\ur{سرنگ} (surang) — a tunnel}

\label{lesson:UR-C316-surang}



\begin{quote}
Before the new one: say the Urdu for a prisoner, then the Urdu for a guard kept.
\end{quote}

\subsection*{You'll want to know: \ur{سرنگ}}
\textbf{\ur{سرنگ}} — \emph{surang} — \textquotedblleft{}a tunnel\textquotedblright{}.

\begin{quote}
\ur{قیدیوں} \ur{نے} \ur{ایک} \ur{سرنگ} \ur{کھودی۔}
\end{quote}

\begin{quote}
\emph{qaidiyoṅ ne ek surang khodī.}
\end{quote}

\begin{quote}
The prisoners dug a tunnel.
\end{quote}

Feminine.

\begin{cousinweb}
\textbf{\ur{سرنگ}} (\emph{surang}) is inherited: CDIAL lists the Hindi cousin \emph{suraṅg}, \textquotedblleft{}a burglar's hole, a mine\textquotedblright{}, under Sanskrit \emph{suruṅgā} (CDIAL 13506). It is feminine: \textbf{\ur{سرنگ} \ur{کھودی}}.
\end{cousinweb}

\begin{grammarlens}[title={What you've built}]
One more word for telling what happened.
\end{grammarlens}

\subsection*{Your turn}
\begin{itemize}
\raggedright
  \item \emph{Say it:} \emph{surang}
  \item \emph{Say it:} \emph{surang}, once more
  \item \emph{Say it:} \emph{qaidiyoṅ ne ek surang khodī.} — The prisoners dug a tunnel.
  \item \emph{Recall:} say the Urdu for a prisoner, then the Urdu for a guard kept, then say \emph{surang} again
\end{itemize}

\begin{tcolorbox}[breakable,colback=teal!4,colframe=teal!35!black,title={Before you move on}]
What does \ur{سرنگ} mean? (\textquotedblleft{}A tunnel\textquotedblright{}.) Say it once more.
\end{tcolorbox}
\section[\texorpdfstring{\ur{کھسکنا} (khisaknā)}{کھسکنا (khisaknā)}]{\ur{کھسکنا} (khisaknā) — to slip away}

\label{lesson:UR-C316-khisakna}



\begin{quote}
Before the new one: say the Urdu for a guard kept, then the Urdu for a tunnel.
\end{quote}

\subsection*{You'll want to know: \ur{کھسکنا}}
\textbf{\ur{کھسکنا}} — \emph{khisaknā} — \textquotedblleft{}to slip away\textquotedblright{}.

\begin{quote}
\ur{قیدی} \ur{رات} \ur{کو} \ur{سرنگ} \ur{سے} \ur{کھسک} \ur{گئے۔}
\end{quote}

\begin{quote}
\emph{qaidī rāt ko surang se khisak gae.}
\end{quote}

\begin{quote}
The prisoners slipped away through the tunnel at night.
\end{quote}

Past (he …): \ur{وہ} \ur{کھسک} \ur{گیا}, \emph{voh khisak gayā}.

\begin{cousinweb}
\textbf{\ur{کھسکنا}} (\emph{khisaknā}) is inherited: CDIAL lists the Hindi cousin \emph{khisaknā}, \textquotedblleft{}to slip away\textquotedblright{}, under a form it reconstructs (CDIAL 3888). It is leaving quietly, so that nobody sees.
\end{cousinweb}

\begin{grammarlens}[title={What you've built}]
One more word for telling what happened.
\end{grammarlens}

\subsection*{Your turn}
\begin{itemize}
\raggedright
  \item \emph{Say it:} \emph{khisaknā}
  \item \emph{Say it:} \emph{khisaknā}, once more
  \item \emph{Say it:} \emph{qaidī rāt ko surang se khisak gae.} — The prisoners slipped away through the tunnel at night.
  \item \emph{Recall:} say the Urdu for a guard kept, then the Urdu for a tunnel, then say \emph{khisaknā} again
\end{itemize}

\begin{tcolorbox}[breakable,colback=teal!4,colframe=teal!35!black,title={Before you move on}]
What does \ur{کھسکنا} mean? (\textquotedblleft{}To slip away\textquotedblright{}.) Say it once more.
\end{tcolorbox}
\section[\texorpdfstring{\ur{رینگنا} (rengnā)}{رینگنا (rengnā)}]{\ur{رینگنا} (rengnā) — to crawl}

\label{lesson:UR-C316-rengna}



\begin{quote}
Before the new one: say the Urdu for a tunnel, then the Urdu for to slip away.
\end{quote}

\subsection*{You'll want to know: \ur{رینگنا}}
\textbf{\ur{رینگنا}} — \emph{rengnā} — \textquotedblleft{}to crawl\textquotedblright{}.

\begin{quote}
\ur{وہ} \ur{سرنگ} \ur{میں} \ur{رینگ} \ur{رہے} \ur{تھے۔}
\end{quote}

\begin{quote}
\emph{voh surang meṅ reng rahe the.}
\end{quote}

\begin{quote}
They were crawling through the tunnel.
\end{quote}

Past (he …): \ur{وہ} \ur{رینگا}, \emph{voh rengā}.

\begin{cousinweb}
\textbf{\ur{رینگنا}} (\emph{rengnā}) is inherited: CDIAL lists the Hindi cousin \emph{rẽgnā}, \textquotedblleft{}to crawl\textquotedblright{}, under Sanskrit \emph{ríṅgati} (CDIAL 10739). Babies, snakes and soldiers under wire all \textbf{\ur{رینگتے} \ur{ہیں}}.
\end{cousinweb}

\begin{grammarlens}[title={What you've built}]
One more word for telling what happened.
\end{grammarlens}

\subsection*{Your turn}
\begin{itemize}
\raggedright
  \item \emph{Say it:} \emph{rengnā}
  \item \emph{Say it:} \emph{rengnā}, once more
  \item \emph{Say it:} \emph{voh surang meṅ reng rahe the.} — They were crawling through the tunnel.
  \item \emph{Recall:} say the Urdu for a tunnel, then the Urdu for to slip away, then say \emph{rengnā} again
\end{itemize}

\begin{tcolorbox}[breakable,colback=teal!4,colframe=teal!35!black,title={Before you move on}]
What does \ur{رینگنا} mean? (\textquotedblleft{}To crawl\textquotedblright{}.) Say it once more.
\end{tcolorbox}
